\section{Input Representation in BERT}
\medskip

Before entering the transformer layers, each token position is represented as the sum of three learned vectors:
\[
\text{Input embedding}
=
E_{\text{token}}
+
E_{\text{position}}
+
E_{\text{segment}}.
\]

Token embeddings map vocabulary items to 768-dimensional vectors. These are static at the input layer. Contextualization occurs only after self-attention layers.

Positional embeddings encode order information. Since attention itself does not know whether a word appears first or last, positional vectors provide this structure.

Segment embeddings indicate whether a token belongs to Sentence~1 or Sentence~2. RoBERTa removes this component.

The summed embeddings are passed through 12 transformer layers in BERT-base. The output is a sequence of contextual token representations:
\[
T_{\text{CLS}}, T_1, \dots, T_K,
\]
each in $\mathbb{R}^{768}$.

\bigskip

